\documentclass[preview]{standalone}

\usepackage{amsmath}
\usepackage{amssymb}
\usepackage{stellar}
\usepackage{definitions}

\begin{document}

\id{i-metodi-della-psicologia}
\genpage

\section{I metodi della psicologia}

\begin{snippet}{metodi-psicologia-overview}
    La psicologia usa principalmente tre metodi:
    \begin{itemize}
        \item il metodo osservativo;
        \item il metodo sperimentale;
        \item il metodo correlazionale.
    \end{itemize}
\end{snippet}

\section{Osservazione}

\begin{snippetdefinition}{osservazione-naturalistica-definition}{Osservazione naturalistica}
    Nell'\emph{osservazione naturalistica}, il ricercatore osserva un
    comportamento che accade in natura, senza cercare di cambiarlo o
    influenzarlo. Il metodo permette di scoprire l'estensione di un fenomeno e
    individuare relazioni significative.
\end{snippetdefinition}

\begin{snippetdefinition}{osservazione-clinica-definition}{Osservazione clinica}
    Nell'\emph{osservazione clinica}, l'osservatore interagisce con la persona
    osservata cercando di ottenere informazioni attraverso domande opportune.
    Il metodo è usato soprattutto a scopi diagnostici e psicoterapeutici.
\end{snippetdefinition}

\section{Metodo sperimentale}

\begin{snippetdefinition}{ambiguita-causale-definition}{Ambiguità causale}
    L'\emph{ambiguità causale} si verifica quando esistono più cause possibili
    per un determinato effetto.
\end{snippetdefinition}

\begin{snippetdefinition}{variabile-psicologia-definition}{Variabile}
    Una \emph{variabile} è qualsiasi caratteristica che può assumere diversi
    valori, quantitativi o qualitativi.
\end{snippetdefinition}

\begin{snippetdefinition}{variabile-indipendente-definition}{Variabile indipendente}
    La \emph{variabile indipendente} è la variabile manipolata dallo
    sperimentatore per verificarne l'effetto su un comportamento o su un altro
    esito misurabile.
\end{snippetdefinition}

\begin{snippetdefinition}{variabile-dipendente-definition}{Variabile dipendente}
    La \emph{variabile dipendente} è la variabile osservata e misurata dallo
    sperimentatore. Il suo valore dovrebbe dipendere dai cambiamenti indotti
    dalla variabile indipendente.
\end{snippetdefinition}

\begin{snippetexample}{colore-stanza-creativita-example}{Colore della stanza e creatività}
    In uno studio sull'effetto del colore della stanza sulla creatività, il
    ricercatore manipola il colore della stanza, per esempio blu o rosso:
    questa è la variabile indipendente. Il punteggio ottenuto dai partecipanti
    in un test di creatività è la variabile dipendente.
\end{snippetexample}

\includesnpt[width=65\%,src=/snippet/static-psychology/independent-dependent-confounding-variables.jpg]{centered-img}

\begin{snippetdefinition}{standardizzazione-procedure-definition}{Standardizzazione delle procedure}
    La \emph{standardizzazione delle procedure} consiste nell'utilizzo di
    procedure uniformi e costanti in ogni fase della sperimentazione e per ogni
    partecipante.
\end{snippetdefinition}

\begin{snippetdefinition}{standardizzazione-norme-definition}{Standardizzazione delle norme}
    La \emph{standardizzazione delle norme} consiste nell'uniformità
    nell'attribuzione dei punteggi e nella valutazione dei dati.
\end{snippetdefinition}

\begin{snippetdefinition}{definizione-operazionale-definition}{Definizione operazionale}
    Una \emph{definizione operazionale} formula un concetto nei termini delle
    procedure utilizzate per determinarlo o misurarlo.
\end{snippetdefinition}

\begin{snippetdefinition}{variabile-confondente-definition}{Variabile confondente}
    Una \emph{variabile confondente} è una variabile non prevista che modifica
    il comportamento dei partecipanti e rende ambigua l'interpretazione dei
    risultati.
\end{snippetdefinition}

\begin{snippetdefinition}{effetto-aspettativa-definition}{Effetto dell'aspettativa}
    L'\emph{effetto dell'aspettativa}, o effetto Rosenthal, consiste nella
    distorsione dei risultati provocata dalle aspettative che lo sperimentatore
    o i soggetti sperimentali hanno riguardo ai risultati.
\end{snippetdefinition}

\begin{snippetdefinition}{effetto-placebo-definition}{Effetto placebo}
    L'\emph{effetto placebo} si verifica quando i partecipanti modificano le
    proprie risposte in assenza di una manipolazione sperimentale efficace.
\end{snippetdefinition}

\begin{snippetdefinition}{procedure-controllo-definition}{Procedure di controllo}
    Le \emph{procedure di controllo} sono procedure finalizzate a mantenere
    costanti tutte le variabili e le condizioni non legate all'ipotesi da
    verificare.
\end{snippetdefinition}

\begin{snippetdefinition}{controllo-singolo-cieco-definition}{Controllo a singolo cieco}
    Il \emph{controllo a singolo cieco} si ha quando i partecipanti non sanno a
    quale condizione sperimentale sono stati assegnati.
\end{snippetdefinition}

\begin{snippetdefinition}{controllo-doppio-cieco-definition}{Controllo a doppio cieco}
    Il \emph{controllo a doppio cieco} si ha quando né i partecipanti né lo
    sperimentatore sanno a quale condizione sperimentale sono stati assegnati i
    soggetti esaminati.
\end{snippetdefinition}

\begin{snippetdefinition}{disegno-ricerca-definition}{Disegno di ricerca}
    Un \emph{disegno di ricerca} è una procedura scientifica che consente
    un'interpretazione non ambigua dei risultati, escludendo a priori
    interpretazioni alternative dovute a possibili variabili confondenti.
\end{snippetdefinition}

\begin{snippetdefinition}{disegno-tra-soggetti-definition}{Disegno tra soggetti}
    Il \emph{disegno tra soggetti} prevede che i partecipanti siano assegnati
    casualmente alla condizione sperimentale o alla condizione di controllo.
\end{snippetdefinition}

\begin{snippetdefinition}{assegnazione-casuale-definition}{Assegnazione casuale}
    L'\emph{assegnazione casuale} è una strategia che attribuisce a ogni
    partecipante la stessa probabilità di trovarsi in una o nell'altra
    condizione sperimentale.
\end{snippetdefinition}

\begin{snippetdefinition}{disegno-entro-soggetti-definition}{Disegno entro i soggetti}
    Il \emph{disegno entro i soggetti} è un disegno sperimentale che utilizza
    ogni partecipante come controllo di se stesso.
\end{snippetdefinition}

\begin{snippetdefinition}{quasi-esperimento-definition}{Quasi esperimento}
    Si parla di \emph{quasi esperimento} quando esiste un'ipotesi causale da
    verificare, è possibile manipolare la variabile indipendente e si possono
    confrontare due o più condizioni sperimentali, ma non sono soddisfatti tutti
    i requisiti di un vero esperimento.
\end{snippetdefinition}

\section{Metodo correlazionale}

\begin{snippetdefinition}{metodo-correlazionale-definition}{Metodo correlazionale}
    Il \emph{metodo correlazionale} è un metodo di ricerca utilizzato per
    stabilire il grado di associazione tra due variabili, senza manipolarle
    sperimentalmente.
\end{snippetdefinition}

\begin{snippetdefinition}{coefficiente-correlazione-definition}{Coefficiente di correlazione}
    Il \emph{coefficiente di correlazione} \(r\) è una misura statistica che
    indica il grado di correlazione esistente tra due variabili. Il suo valore
    è compreso tra \(+1{,}0\) e \(-1{,}0\).
\end{snippetdefinition}

\begin{snippetdefinition}{correlazione-positiva-definition}{Correlazione positiva}
    Una \emph{correlazione positiva} indica che, all'aumentare del valore di
    una variabile, aumenta anche il valore dell'altra.
\end{snippetdefinition}

\begin{snippetdefinition}{correlazione-negativa-definition}{Correlazione negativa}
    Una \emph{correlazione negativa} indica che, all'aumentare del valore di
    una variabile, il valore dell'altra diminuisce.
\end{snippetdefinition}

\begin{snippetnote}{correlazione-causalita-note}{Correlazione e causalità}
    La correlazione non implica causalità. Un'associazione tra due variabili
    può dipendere da molteplici nessi causa-effetto o da una terza variabile
    sottostante.
\end{snippetnote}

\end{document}
